\exercisesection{星代数}

\begin{enumerate}
\def\labelenumi{\arabic{enumi})}
\item
  找出一个幺对合 Banach 代数 A 和一个酉元 \(u \in A\)，使 \(\|u\| \neq 1\)。
\item
  设 A 为 Banach 代数，\(A_{1} \subset A\) 为其闭子代数。若 \(A_{1}\) 是对合代数，则 \(A_{1}\) 的每个特征标都可以延拓为 A 的特征标。
\item
  设 A 为幺对合代数。假设对 A 中的每个 x，x 的谱都不含 −1。
\end{enumerate}

\begin{enumerate}
\def\labelenumi{\alph{enumi})}
\item
  对 A 的每个幂等元 x，证明存在 A 的一个厄米幂等元 y，使 \(xA = yA\)。（考虑 \(1 + (x - x^{*})(x^{*} - x)\) 的逆 z；证明 z 与 x 交换，并令 \(y = xx^{*}z\)。）
\item
  对 A 的每个幂等元 \(x\)，证明存在 A 的一个厄米幂等元 \(y\) 与 A 中可逆的 \(u\)，使 \(x = uyu^{-1}\)。（若 \(x\mathrm{A} = y\mathrm{A}\)，则令 \(u = (1 - (y - x))^{-1}\)。）
\end{enumerate}

\begin{enumerate}
\def\labelenumi{\arabic{enumi})}
\setcounter{enumi}{3}
\item
  设 A 为赋范代数，带有一个连续对合。若对一切 \(\|x\|' = \sup(\|x\|, \|x^{*}\|)\) 令 \(x \in A\)，则 A 成为一个对合赋范代数，且新范数与旧范数等价。
\item
  设 A 为无根交换 Banach 代数。A 的每个对合都是连续的。（利用 I，第 40 页的命题 9。）
\end{enumerate}

\begin{enumerate}
\def\labelenumi{¶ \arabic{enumi})}
\setcounter{enumi}{5}
\tightlist
\item
  设 A 为每个特征标都是厄米的交换幺对合 Banach 代数。假设对每个不可逆的厄米元 \(x \in A\)，有 \(\|(x - \lambda)^{-1}\| = o(\mathcal{I}(\lambda)^{-2})\)（当 \(\lambda\) 保持 \(\mathcal{I}(\lambda) \neq 0\) 而趋于 0 时）。则 A 的每个闭理想 I 都是一些极大理想的交。（设 \(\varphi: A \longrightarrow A/I\) 为典范态射。设 \(x \in A\) 使 \(\varphi(x)\) 属于 A/I 的根。要证明 \(\varphi(x) = 0\)。化归到 \(x\) 是厄米的情形。然后利用 I，第 161 页的习题 29。）
\end{enumerate}

\begin{enumerate}
\def\labelenumi{\arabic{enumi})}
\setcounter{enumi}{6}
\item
  设 A 为交换幺对合 Banach 代数。为使 A 的每个特征标都是厄米的，充要条件是对每个 \(1 + xx^{*}\)，\(x \in A\) 可逆。（若 A 的每个特征标都厄米，则 \(\mathcal{G}(1 + xx^{*})\) 在 \textgreater{} 0（在 \(\mathsf{X}(\mathsf{A})\) 上）处处成立。若 A 的特征标 \(\chi\) 非厄米，则存在厄米元 \(x \in A\) 使 \(\chi(x) = i\)，从而 \(\chi(1 + xx^{*}) = 0\)，而 \(1 + xx^{*}\) 不可逆。）
\item
  构造一个群 G 的例子，使对合 Banach 代数 \(\mathcal{M}^{1}(G)\)（I，第 100 页，例 4）不是星代数。
\item
  设 A 为幺 Banach 代数。A 的酉元所成的群是 A* 的一个极大有界子群。
\item
  设 A 为对合 Banach 代数，\(\widetilde{\mathsf{A}}\) 为添加单位元后所得的代数。证明：I，第 15 页上在 \(\widetilde{\mathsf{A}}\) 上定义的范数，一般说来与 I，第 105 页命题 3 所定义的范数不相一致。
\item
  设 A 为对合 Banach 代数，其对合满足：对一切 \(\|x^{*}x\| = \|x\| \cdot \|x^{*}\|\) 有 \(x \in A\)。则 A 是一个星代数。（对厄米元 y 有 \(\|y^{2}\| = \|y\|^{2}\)，从而 \(\|y\| = \varrho(y)\)；注意到对一切 \(\mathrm{Sp}_{\mathrm{A}}^{\prime}(x^{*}) = \overline{\mathrm{Sp}_{\mathrm{A}}^{\prime}(x)}\) 有 \(x \in A\)，因此我们有
\end{enumerate}

\[
\| x ^ {*} \| \cdot \| x \| = \| x ^ {*} x \| = \varrho (x ^ {*} x) \leqslant \varrho (x ^ {*}) \varrho (x) = \varrho (x ^ {2}) \leqslant \| x \| ^ {2},
\]

从而 \(\|x^{*}\|\leqslant\|x\|\) 且 \(\|x^{*}\|=\|x\|\)。

\begin{enumerate}
\def\labelenumi{\arabic{enumi})}
\setcounter{enumi}{11}
\item
  设 A 为星代数，N 为 A 的正规元所成的集合，\(x \in N\)，V 为 C 中 0 的邻域。存在 x 在 N 中的一个邻域 U，使得对每个 \(y \in U\)，有 \(\mathrm{Sp}_{\mathrm{A}}(y) \subset \mathrm{Sp}_{\mathrm{A}}(x) + \mathrm{V}\) 与 \(\mathrm{Sp}_{\mathrm{A}}(x) \subset \mathrm{Sp}_{\mathrm{A}}(y) + \mathrm{V}\)。（利用 I，第 108 页命题 1 的推论、I，第 23 页的命题 3 以及 I，第 76 页的命题 10。）
\item
  设 A 为 C 上的幺 Banach 代数，\(x \in A\)，且 \(S = \mathrm{Sp}_{\mathrm{A}}(x)\)。假设存在从 \(\varphi\) 到 A 中的一个态射 \(\mathcal{C}(\mathrm{S})\)，它把 1 映为 1、把 z 映为 x（z 表示 S 上的恒等映射）。则对每个 \(\|\varphi(f)\| \geqslant \|f\|\)，有 \(f \in \mathcal{C}(\mathrm{S})\)。特别地，若态射 \(\varphi\) 连续，则 \(\varphi\) 是到其像上的一个同胚。（可以假设 A 交换。设 \(\lambda \in S\)。存在 \(\chi \in \mathrm{X}(\mathrm{A})\) 使 \(\lambda = \chi(x) = \chi(\varphi(z)) = (\mathrm{X}(\varphi)(\chi))(z)\)。于是 \(\mathrm{X}(\varphi)(\chi)\) 就是 \(f \mapsto f(\lambda)\) 的特征标 \(\mathcal{C}(\mathrm{S})\)。设 \(f \in \mathcal{C}(\mathrm{S})\)。由前所证，存在 \(\chi \in \mathrm{X}(\mathrm{A})\) 使 \(|f|\) 在 \(\mathrm{X}(\varphi)(\chi)\) 处达到其最大值。于是 \(|\chi(\varphi(f))| = \|f\|\)，从而 \(\|\varphi(f)\| \geqslant \|f\|\)。）
\item
  设 G 为局部紧群。若 Stell(G) 有单位元，则群 G 是离散的。（利用下述事实：在 Hilbert 空间 \(L^{2}(G)\) 中，恒等映射是诸自同态 \(\gamma_{f}\) 的范数极限，其中 \(\gamma\) 是左正则表示，而 \(f \in L^{1}(G)\)。）
\item
  \begin{enumerate}
  \def\labelenumii{\alph{enumii})}
  \tightlist
  \item
    设 A 为交换星代数。若 x 与 y 是 A 中满足 \(x \leqslant y\) 的正元素，则 \(x^{2} \leqslant y^{2}\)。
  \end{enumerate}
\end{enumerate}

\begin{enumerate}
\def\labelenumi{\alph{enumi})}
\setcounter{enumi}{1}
\item
  存在一个星代数 A 及 A 中元素 x、y，使得 \(0 \leqslant x \leqslant y\)，但 \(x^{2}\) 并不成立\(\leqslant y^{2}\)。
\item
  存在一个星代数 A 及 A 中元素 \(x\) 与 \(y\)，使得 \(|x + y|\) 并不小于等于 \(\leqslant |x| + |y|\)。
\end{enumerate}

\begin{enumerate}
\def\labelenumi{\arabic{enumi})}
\setcounter{enumi}{15}
\tightlist
\item
  设 X 为拓扑空间。设 \(\mathcal{C}_{b}(\mathrm{X})\) 为 X 上有界连续复值函数所成的交换星代数。
\end{enumerate}

\begin{enumerate}
\def\labelenumi{\alph{enumi})}
\item
  设 \(\widehat{\mathrm{X}} = \mathrm{X}(\mathcal{C}_b(\mathrm{X}))\)；空间 \(\widehat{\mathrm{X}}\) 是一个紧空间，使得 \(\mathcal{C}(\widehat{\mathrm{X}})\) 典范地等同于 \(\mathcal{C}_b(\mathrm{X})\)。
\item
  映射 \(j: X \to \widehat{X}\)（把 x 对应到特征标 \(f \mapsto f(x)\)）是连续的。
\item
  若 X 是完全正则的（TG，IX，第 8 页，定义 4），则 j 是 X 到 \(\widehat{X}\) 的一个稠密子空间上的同胚。（为证 j 的像在 \(\widehat{X}\) 中稠密，证明：在 \(\widehat{X}\) 上连续且在 \(j(\mathrm{X})\) 上为零的复值函数恒为零。）
\item
  空间 \(\widehat{X}\) 典范地等同于 X 的 Stone–Čech 紧化（TG，IX，第 9 页）。
\end{enumerate}

\begin{enumerate}
\def\labelenumi{\arabic{enumi})}
\setcounter{enumi}{16}
\item
  设 A 为 C*-代数，\(x \in \mathrm{A}\)。若 \(x = a - b\)（其中 \(a, b\) 为正元素且 \(ab = ba = 0\)），则 \(a = x^{+}\) 且 \(b = x^{-}\)。
\item
  设 A 为可分星代数。存在 \((e_{n})_{n\in\mathbf{N}}\) 中的一个序列 \(A_{+}\)，使相联的基本滤子（TG，I，第 42 页，定义 7）是 A 的一个递增近似单位元。
\item
  设 A 为星代数。称元素 \(x\) 是系统正的，如果 \(x \in \mathrm{A}_{+}\) 且 \(x\mathrm{A}x\) 在 A 中稠密。
\end{enumerate}

\begin{enumerate}
\def\labelenumi{\alph{enumi})}
\item
  证明：交换星代数 A 含有系统正元素，充要条件是存在一个在无穷远处可数的局部紧拓扑空间 X，使 A 同构于 \(\mathcal{C}_{0}(X)\)。
\item
  若 A 是幺的，则 A 的元素 x 是严格正的，当且仅当 x 是正的且可逆。
\item
  设 \(x \in A\) 是一个系统正元素，满足 \(\|x\| \leqslant 1\)。证明序列 \((x^{1/n})_{n \geqslant 1}\) 是 A 中的一个递增近似单位元。
\item
  星代数 A 含有系统正元素，充要条件是存在 A 中的一个序列 \((x_{n})_{n\geqslant1}\)，它构成 A 的一个近似恒等元。
\end{enumerate}

\begin{enumerate}
\def\labelenumi{\arabic{enumi})}
\setcounter{enumi}{19}
\item
  设 A 为 C*-代数，I 为 A 的一个闭双边理想，B 为 A 的一个闭对合子代数。则 B+I 是 A 的一个闭对合子代数，且存在等距同构 \(B + I/I \rightarrow B/B \cap I\)。
\item
  设 S 为 C 中的闭单位圆盘。设 A 为从 S 到 C 的、在开单位圆盘中全纯的连续函数所成的对合 Banach 代数，其对合由 \(f^{*}(z) = \overline{f(\overline{z})}\) 给出。A 中存在一个厄米元素，其谱不含于 R 中。
\item
  设 \(\mathrm{A} = \mathcal{C}([0,1])\)。证明由 {[}0, 1{]} 上的恒等函数所生成的 A 的理想 I 不是闭的。
\item
  设 A 为幺 C*-代数。用 \(A_{+}^{\prime}\) 表示 A 上的线性形式 f（不必连续）所成的向量空间，使得
\end{enumerate}

\[
f (\mathrm{A} _ {+}) \subset [ 0, + \infty [
\]

（“正线性形式”）。

\begin{enumerate}
\def\labelenumi{\alph{enumi})}
\item
  设 \(f \in A_{+}^{\prime}\)。对 \((x_{n})_{n \in \mathbf{N}}\) 中的每个序列 \(A_{+}^{\leqslant 1}\)，族 \((f(x_{n}))\) 有界。（证明：对 \(\sum_{i} f(x_{i})y_{i}\) 中的每个正族 \((y_{i})\)，级数 \(\ell^{1}(\mathbf{N})\) 收敛。）
\item
  线性映射 f 是连续的。
\item
  线性形式 \(f\) 属于 \(\mathrm{A}_{+}^{\prime}\) 的充要条件是 \(f\) 连续且 \(\| f \| = f(1)\)。
\item
  设 \(x \in A_{h}\) 为厄米元素。存在正线性形式 \(f \in A_{+}^{\prime}\)，使 \(|f(x)| = \|x\|\)。（利用 I，第 174 页的习题 7。）
\end{enumerate}

\begin{enumerate}
\def\labelenumi{\arabic{enumi})}
\setcounter{enumi}{23}
\tightlist
\item
  设 A 为 C 上的幺 Banach 代数。假设 \(x \mapsto x^{*}\) 与 \(x \mapsto x^{\circ}\) 是 A 上的对合，且 A 赋以其中任一对合都成为星代数。用 \(A_{1}\) 与 \(A_{2}\) 记这两个星代数。
\end{enumerate}

\begin{enumerate}
\def\labelenumi{\alph{enumi})}
\item
  沿用习题 23 的记号，证明 \(\mathrm{A}_{1,+}^{\prime} = \mathrm{A}_{2,+}^{\prime}\)。用 \(\mathrm{A}_{+}^{\prime}\) 记这个集合。
\item
  设 \(x \in A_{1,h}\) 是满足 \(x^{*} = x\) 的元素。证明：对每个 \(f \in A_{+}^{\prime}\)，有 \(f(i(x - x^{\circ})) = 0\)。由此推出 \(x \in A_{2,h}\)。（利用习题 23 的 b)。）
\item
  由此得出结论：\(A_{1} = A_{2}\)，并且一个复 Banach 代数至多有一个使它成为 C*-代数的对合。
\end{enumerate}

\begin{enumerate}
\def\labelenumi{\arabic{enumi})}
\setcounter{enumi}{24}
\tightlist
\item
  设 G 为群。用 C{[}G{]} 表示 G 在 C 上的群代数。它是 Hilbert 空间 \(\ell^{2}(\mathrm{G})\) 的一个子空间。对 G 中的 g，用 {[}g{]} 表示 C{[}G{]} 中对应于 g 的元素。诸 {[}g{]}（对 \(g \in G\)）构成 \(\ell^{2}(G)\) 的一个标准正交基，且 C{[}G{]} 在 \(\ell^{2}(G)\) 中稠密。
\end{enumerate}

\begin{enumerate}
\def\labelenumi{\alph{enumi})}
\tightlist
\item
  映射
\end{enumerate}

\[
\sum_ {g} \alpha_ {g} [ g ] \mapsto \sum_ {g} \overline {{\alpha}} _ {g} [ g ^ {- 1} ]
\]

是 \(\mathbf{C}[\mathrm{G}]\) 上的一个对合。

\begin{enumerate}
\def\labelenumi{\alph{enumi})}
\setcounter{enumi}{1}
\tightlist
\item
  对 \(g \in G\)，设 \(\pi(g)\) 为 C\(x \mapsto gx\)[{G}]{ 的映射 }。把 g 对应到 \(\pi(g)\) 的映射可延拓为从 C{[}G{]} 到 \(\mathcal{L}(\ell^{2}(G))\) 中的一个连续单射对合代数态射。
\end{enumerate}

用 \(\operatorname{Stell}_{r}(\mathbf{G})\) 记在 \(\mathcal{L}(\ell^{2}(\mathbf{G}))\) 中映射 \(\pi\) 的像的闭包，称之为 G 的约化 C*-代数。它是一个 C*-代数。我们把 \(\mathbf{C}[\mathbf{G}]\) 经 \(\pi\) 在 \(\operatorname{Stell}_{r}(\mathbf{G})\) 中的像与它等同。

\begin{enumerate}
\def\labelenumi{\alph{enumi})}
\setcounter{enumi}{2}
\tightlist
\item
  线性迹映射 \(\operatorname{Tr} \colon \mathbf{C}[\mathrm{G}] \to \mathbf{C}\)（满足 \(\operatorname{Tr}(1) = 1\) 且对 \(\operatorname{Tr}(g) = 0\) 有 \(g \neq 1\)）可延拓为连续线性映射 \(\operatorname{Stell}_r(\mathrm{G}) \to \mathbf{C}\)。
\end{enumerate}

\begin{enumerate}
\def\labelenumi{\alph{enumi})}
\setcounter{enumi}{3}
\tightlist
\item
  我们有 \(\operatorname{Tr}(xy) = \operatorname{Tr}(yx)\)（对一切 \(x\) 与 \(y\)，它们都属于 \(\operatorname{Stell}_r(G) \)）。
\end{enumerate}

\begin{enumerate}
\def\labelenumi{\alph{enumi})}
\setcounter{enumi}{4}
\item
  我们有：对一切 \(\operatorname{Tr}(xx^{*}) \geqslant 0\) 属于 \(x\)，\(\operatorname{Stell}_r(G)\)；且 \(\operatorname{Tr}(xx^{*}) = 0\) 当且仅当 \(x = 0\)。
\item
  设 e 是 C{[}G{]} 的一个幂等元，满足 \(e \neq 0\) 且 \(e \neq 1\)。证明 \(0 < \operatorname{Tr}(e) < 1\)。（利用习题 3。）
\item
  设 e 是 Z{[}G{]} 的一个幂等元。证明 e = 0 或 e = 1。（“Kaplansky 定理”。）
\item
  设 x 是 C{[}G{]} 的一个元素。则 x 可逆，当且仅当它左可逆，当且仅当它右可逆。
\item
  设 n 为严格正的整数。类型为 \((n, n)\) 且系数在 C{[}G{]} 中的矩阵 x 可逆，当且仅当它左可逆，当且仅当它右可逆。
\end{enumerate}

\begin{enumerate}
\def\labelenumi{\arabic{enumi})}
\setcounter{enumi}{25}
\item
  设 A 为幺 C*-代数。设 \(a\) 与 \(b\) 为 A 的正规元，\(x \in \mathrm{A}^*\) 满足 \(b = xax^{-1}\)。存在 A 的酉元 \(u\) 使 \(b = uau^*\)。（利用前一习题以及 \(x\) 的极分解，注意 \(|x|\) 属于 \(x\) 的双交换子之中。）
\item
  设 A 为幺 C*-代数，a、b、c 为 A 的元素。假设 b 与 c 正规。若 ac = cb，则对 a 的谱与 b 的谱的并上的每个连续函数 f，有 \(f(a)c = cf(b)\)。
\item
  设 A 为 C*-代数。设 \(a\) 与 \(b\) 为 A 的正元素，满足 \(ab \neq 0\)。则 \(\mathrm{Sp}_{\mathrm{A}}(ab)\) 不化为 0。
\item
  \begin{enumerate}
  \def\labelenumii{\alph{enumii})}
  \tightlist
  \item
  设 G 为局部紧拓扑群。Banach 代数 L\(^{1}\)(G) 具有一个近似恒等元。（利用 INT，VIII，§2，第 7 节，引理 4。）
  \end{enumerate}
\end{enumerate}

\begin{enumerate}
\def\labelenumi{\alph{enumi})}
\setcounter{enumi}{1}
\item
  设 A 为具有近似恒等元的 Banach 代数。对每个特征标 \(\chi\in\mathsf{X}(\mathrm{A})\)，有 \(\|\chi\|=1\)。
\item
  L\(^{1}\)(Z) 的由序列（\(x_{n}\)）——满足 \(x_{n} = 0\)（当 \(n < 0\)）——所组成的 Banach 子代数没有近似恒等元。（构造一个范数 \(< 1\) 的非零特征标。）
\end{enumerate}

\begin{enumerate}
\def\labelenumi{\alph{enumi})}
\setcounter{enumi}{3}
\tightlist
\item
  设 G 为由两个元素 \(a\) 与 \(b\) 生成的自由群，\(\mu\) 为 G 上的一个 Haar 测度。设 I 为 \(\mathrm{L}^1 (\mathrm{G})\) 中由使 \(f\in \mathrm{L}^{1}(\mathrm{G})\) 的 \(\mu (f) = 0\) 所成的闭理想。令 \(\xi = e^{2i\pi /3}\)。令 \(x = \varepsilon_{e} + \xi \varepsilon_{a} + \xi^{-1}\varepsilon_{b}\in \mathrm{L}^{1}(\mathrm{G})\)。对每个 \(f\in \mathrm{L}^{1}(\mathrm{G})\)，有 \(\| x*f - x\| \geqslant 1\)。Banach 代数 I 没有近似恒等元。
\end{enumerate}

\begin{enumerate}
\def\labelenumi{¶ \arabic{enumi})}
\setcounter{enumi}{29}
\tightlist
\item
  设 A 为具有近似恒等元的 Banach 代数。用 \(\widetilde{\mathsf{A}}\) 表示由 A 添加单位元（记作 1）所得的 Banach 代数。
\end{enumerate}

\begin{enumerate}
\def\labelenumi{\alph{enumi})}
\item
  设 \(x \in \mathrm{A}\)。则 \(x\) 属于由 \(x\) 生成的闭左理想。
\item
  对每个 \(\varepsilon > 0\) 与 \((x_i)_{i \in \mathrm{I}}\) 中元素的每个有限族 \(\mathrm{A}\)，存在 \(e \in \mathrm{A}\) 使 \(\| e \| \leqslant 1\)，且对一切 \(\| ex_i - x_i \| < \varepsilon\) 有 \(i \in \mathrm{I}\)。
\item
  设 \(x \in A\)、\(\varepsilon > 0\) 与 \(0 < \gamma \leqslant 1/8\)。存在 \(\eta > 0\)，使得对满足 \(e \in A\) 与 \(\|e\| \leqslant 1\) 的每个 \(\|ex - x\| < \eta\)，元素 \(\alpha = (1 - \gamma) \cdot 1 + \gamma e\) 在 \(\widetilde{A}\) 中可逆，且 \(\|x - \alpha^{-1}x\| < \varepsilon\)。
\end{enumerate}

\begin{enumerate}
\def\labelenumi{\alph{enumi})}
\setcounter{enumi}{3}
\tightlist
\item
  设 \(x \in \mathrm{A}\)、\(\varepsilon > 0\)、\(\delta > 0\) 与 \(0 < \gamma \leqslant 1/8\)。存在 A 中满足下列性质的序列 \((e_n)_{n \geqslant 1}\)：

\begin{enumerate}
\def\labelenumi{(\roman{enumi})}
\item
  对一切 \(\|e_{n}\| < 2\)，有 \(n \geqslant 1\)；
\item
  元素
\end{enumerate}

\[
y _ {n} = \sum_ {k = 1} ^ {n} \gamma (1 - \gamma) ^ {k - 1} e _ {k} + (1 - \gamma) ^ {n} \cdot 1
\]

在 \(\widetilde{\mathrm{A}}\) 中（对一切 \(n\geqslant 1\)）可逆，且

\[
\left\| y _ {n} ^ {- 1} x - y _ {n - 1} ^ {- 1} x \right\| <   \delta 2 ^ {- n} \quad (n \geqslant 2), \quad \left\| y _ {1} ^ {- 1} x - x \right\| <   \delta / 2.
\]

（假设 \((e_{1},\ldots,e_{n})\) 已经定义，考虑元素 \(e_{n+1}\)（任意，范数 \textless{} 2）；注意

\[
y _ {n + 1} = ((1 - \gamma) \cdot 1 + e _ {n + 1}) \left(\sum_ {k = 1} ^ {n} \gamma (1 - \gamma) ^ {k - 1} e _ {k} ^ {\prime} + (1 - \gamma) ^ {n} \cdot 1\right),
\]

其中 \(e_k' = ((1 - \gamma) \cdot 1 + e_{n+1})^{-1} e_k\)，并确定满足所需条件的 \(e_{n+1}\)（利用 \(\widetilde{\mathsf{A}}\) 的可逆元所成集合是开集、且映射 \(a \mapsto a^{-1}\) 连续这些事实）。
\end{enumerate}

\begin{enumerate}
\def\labelenumi{\alph{enumi})}
\setcounter{enumi}{4}
\tightlist
\item
  对每个 \(x \in A\)，存在 A 中的 y 与 z 使 x = yz（“Cohen 因式分解定理”）。设 \(\delta > 0\)。可以假设 z 属于由 \(x\) 生成的闭左理想，且 \(\|y-z\|<\delta\)。（考虑如上的序列 \((e_{n})_{n\geqslant1}\)，并按之前的方式定义 \(y_{n}\)；序列 \(z_{n}=y_{n}^{-1}x\) 收敛于一个元素 \(z\in\mathrm{A}\)，且 \(z_{n}\) 属于由 \(x\) 生成的闭左理想；序列 \((y_{n})\) 收敛于 \(y\in\mathrm{A}\)，且 \(x=yz\)。）
\end{enumerate}

\begin{enumerate}
\def\labelenumi{\arabic{enumi})}
\setcounter{enumi}{30}
\tightlist
\item
  设 A 为关于 Lebesgue 测度在 {[}0, 1{]} 上可积的复函数所成的 Banach 空间。令
\end{enumerate}

\[
(f \cdot g) (t) = \int_ {0} ^ {t} f (t - x) g (x) d x
\]

对 \(t \in [0, 1]\)。

\begin{enumerate}
\def\labelenumi{\alph{enumi})}
\item
  函数 \(f \cdot g\) 几乎处处有定义，且是 A 的一个元素。
\item
  赋以运算 \((f,g)\mapsto f\cdot g\) 后，空间 A 是一个交换 Banach 代数。
\item
  设 \(x \in \mathrm{A}\) 是恒等于 1 的常值函数。函数 \(x^n\) 是
\end{enumerate}

\[
t \mapsto \frac {t ^ {n - 1}}{(n - 1) !}.
\]

\begin{enumerate}
\def\labelenumi{\alph{enumi})}
\setcounter{enumi}{3}
\item
  代数 A 由 x 拓扑地生成，元素 x 是拟幂零元，且 A 等于它的根。
\item
  代数 A 具有近似恒等元。
\item
  代数 A 没有极大理想，无论是正则的还是非正则的，无论闭与否（利用前一习题以及 I，第 166 页的习题 1）。特别地，X(A) 是空集，因而是紧的。
\end{enumerate}

\begin{enumerate}
\def\labelenumi{\arabic{enumi})}
\setcounter{enumi}{31}
\tightlist
\item
  设 G 为群。称 G 具有 Powers 性质，如果对每个 \(n \in N\) 以及 G = \(G = \{e\}\) 的每个有限子集 F，存在 G 分成两个子集 A 与 B 的一个划分，以及 G 的元素的一个族 \((g_i)_{1 \leqslant i \leqslant n}\)，使得
\end{enumerate}

\begin{enumerate}
\def\labelenumi{(\roman{enumi})}
\item
  对一切 \(g \in \mathrm{F}\)，集合 A 与 gA 不相交；
\item
  对任意两个不同的整数 \(j\) 与 \(k\)（介于 1 与 \(n\) 之间），集合 \(g_{j}\mathrm{B}\) 与 \(g_{k}\mathrm{B}\) 不相交。
\end{enumerate}

设 G 为具有 Powers 性质的群。设 I 为约化 C*-代数 \(\text{Stell}_{r}(\text{G})\)（习题 25）的一个非零闭双边理想。设 \(\pi: C[G] \to \text{Stell}_{r}(G)\) 为典范表示。

\begin{enumerate}
\def\labelenumi{\alph{enumi})}
\tightlist
\item
  存在一个族 \((\lambda_g)_{g\in \mathrm{G - }\{e\}}\)，其值属于 \(\mathbf{C}\)，使得
\end{enumerate}

\[
x = \sum_ {g \in \mathrm{G} - \{e \}} \lambda_ {g} \pi (g)
\]

是 \(\operatorname{Stell}_{r}(\mathrm{G})\) 的一个厄米元素，且 \(y = 1 + x\) 属于 I。

\begin{enumerate}
\def\labelenumi{\alph{enumi})}
\setcounter{enumi}{1}
\tightlist
\item
  设 \(\varepsilon > 0\) 满足 \(\varepsilon < 1\)。存在 G = \{e\} 的一个有限子集 F 以及一个整数 \(n \geqslant 1\)，使得元素
\end{enumerate}

\[
x ^ {\prime} = \sum_ {g \in \mathrm{F}} \lambda_ {g} \pi (g)
\]

满足 \(\|x' - x\| < \varepsilon\) 与 \(\|x'\| < \frac{\sqrt{n}}{2}(1 - \varepsilon)\)。

\begin{enumerate}
\def\labelenumi{\alph{enumi})}
\setcounter{enumi}{2}
\tightlist
\item
  取 \(G = A \cup B\) 与 \((g_{i})_{1 \leqslant i \leqslant n}\)，它们对有限子集 F 与整数 n 满足 Powers 性质的条件。令
\end{enumerate}

\[
y ^ {\prime} = \frac {1}{n} \sum_ {i = 1} ^ {n} \pi (g _ {i}) x ^ {\prime} \pi (g _ {i} ^ {- 1}).
\]

我们有 \(\|y'\| < 1 - \varepsilon\)。（设 \(p_i\) 是 \(\ell^{2}(\mathrm{G})\) 到 \(\ell^{2}(g_{i}\mathrm{B})\) 上的正交投影；验证

\[
(1 - p _ {i}) \pi (g _ {i}) x ^ {\prime} \pi (g _ {i} ^ {- 1}) (1 - p _ {i}) = 0
\]

对 \(1 \leqslant i \leqslant n\) 成立，并由此推出

\[
1 + y ^ {\prime} = 1 + \frac {1}{n} \sum_ {i = 1} ^ {n} p _ {i} \pi (g _ {i}) x ^ {\prime} \pi (g _ {i} ^ {- 1}) + \left(\frac {1}{n} \sum_ {i = 1} ^ {n} p _ {i} \pi (g _ {i}) x ^ {\prime} \pi (g _ {i} ^ {- 1}) (1 - p _ {i})\right) ^ {*},
\]

然后注意自同态 \(p_{i}\pi(g_{i})x^{\prime}\pi(g_{i}^{-1})\) 具有两两正交的像，以此估计它们的范数。

\begin{enumerate}
\def\labelenumi{\alph{enumi})}
\setcounter{enumi}{3}
\tightlist
\item
  理想 I 包含元素
\end{enumerate}

\[
1 + \frac {1}{n} \sum_ {i = 1} ^ {n} \pi (g _ {i}) x \pi (g _ {i} ^ {- 1}),
\]

并且这个元素可逆。由此推出，星代数 \(\text{Stell}_{r}(\text{G})\) 除 \(\{0\}\) 与 \(\text{Stell}_{r}(\text{G})\) 外不含任何闭双边理想。

\begin{enumerate}
\def\labelenumi{\alph{enumi})}
\setcounter{enumi}{4}
\tightlist
\item
  设 G 为非交换自由群。则 G 具有 Powers 性质。
\end{enumerate}
% semantic-fragments: legacy-input-seam
\relax{}
